
This paper establishes that scaling and sign-flip symmetry orbits are geometrically large in the Schürholt MNIST MLP zoo --- every oracle-constructed orbit pair across 2,500 measurements exceeds the 0.05 cosine distance significance threshold, with mean values of 0.32 (scaling) and 1.07 (sign-flip). The Neural Functional Transformer is measurably non-invariant to scaling orbits (gap=+0.0238, CI=[+0.0232, +0.0245]) and measurably invariant to sign-flip orbits by construction (gap=$-0.0007$, CI=[$-0.0008, -0.0005])$. This asymmetry --- not designed but emergent from row-level weight tokenization --- implies that scaling canonicalization is the actionable preprocessing target for NFT-family encoders. The majority-sign sign-flip algorithm produces a unique canonical form for only 14.4\% of Schürholt MNIST zoo models; the 85.6\% failure rate is accurately predicted by binomial combinatorics for even d\_in=784, providing a novel structural characterization of the algorithm's limitations.

Property prediction experiments at N=500 are severely underpowered: all Spearman $\rho$ values are near zero and all condition comparisons have CI width $\approx 0.6$. Point estimates for $\Delta\rho_{D}$-A are in the predicted direction (+0.058, +0.053, +0.077 across three properties), but no comparison achieves statistical significance. An unexpected result --- Condition E (random normalization) consistently exceeding Condition D (full canonicalization) --- is most plausibly explained by the non-unique sign-flip canonicalization identified in the audit, though this attribution cannot be confirmed at N=500. Adequate evaluation of the property prediction improvement claim requires $N\geq 5,000$ models and a corrected sign-flip implementation (odd d\_in or magnitude-based tie-breaking).

Together, these findings transform the argument for symmetry canonicalization in weight space learning from a theoretical expectation into a measurement-grounded framework: orbit characterization establishes which symmetries are geometrically large; encoder probing establishes which a given encoder already handles; and algorithm auditing establishes which standard preprocessing approaches fail and why. The most direct next step --- scaling-only canonicalization at full zoo scale --- is specified precisely enough to be immediately implemented.

